\chapter{Conclusão}

Este capítulo sintetiza as principais contribuições deste trabalho de conclusão de curso, revisitando o alcance dos objetivos propostos, destacando as limitações identificadas durante o desenvolvimento e apontando oportunidades para trabalhos futuros.

\section{Síntese do Trabalho e Contribuições}

\textcolor{red}{(Apresentar aqui o resumo completo do projeto: demonstrar como a integração entre Engenharia de Software, microsserviços conteinerizados em Java/Spring Boot, Arquitetura Orientada a Eventos (EDA) e práticas de Spec-Kit/TDD permitiu transformar um processo fabril frágil e baseado em planilhas manuais em uma plataforma escalável, resiliente e totalmente auditável para o ciclo metrológico industrial.)}

\section{Trabalhos Futuros}

\textcolor{red}{(Indicar aqui as sugestões de continuidade e evolução do sistema, tais como: a orquestração avançada dos contêineres com Kubernetes em ambientes multi-cloud distribuídos, o desenvolvimento de modelos preditivos baseados em séries temporais/Machine Learning para antecipar desvios e derivas metrológicas em equipamentos críticos, e a integração com leitores RFID/QR Code no chão de fábrica para localização física de instrumentos.)}
